\origpage{230}{ROBERTSON AND SEYMOUR}
The paper falls into two parts. In the first part (Sections 2--5) we prove a generalization of Kruskal's theorem (1.2); and apply it to prove that, for every containment relation defined on a class of rooted hypergraphs satisfying certain axioms, if it forms a well-quasi-order (in a sense) on the hypergraphs in the class of bounded size, then it forms a well-quasi-order on the hypergraphs in the class of bounded tree-width. (1.5) is a consequence of this where our hypergraphs are graphs. In the second part of the paper (Sections 6--9) we discuss a concrete (but rather complex) containment relation on a class of hypergraphs with additional structure, so-called \textit{patchworks}. We verify that this relation does indeed satisfy our axioms for a containment relation, and so the results of the first part apply to it. These results, applied to patchworks, will be important in future papers.

\sectionhead{2. A Lemma about Trees}

In this section we prove a lemma about rooted trees generalizing (1.2). Most of the trees we will need in this paper are ``rooted'' trees, and so for convenience we define an \textit{undirected tree} to be a non-null connected finite undirected graph without circuits, and a \textit{tree} $T$ to be a directed graph, the undirected graph underlying which is an undirected tree, such that every vertex of $T$ is the head of at most one edge. It follows that there is a unique vertex of each tree $T$ (called the \textit{root} and denoted by $o(T)$) which is the head of no edge of $T$, and every edge of $T$ is directed away from the root.

We begin with a preliminary form of our lemma. If $M$ is a (possibly infinite) graph, its vertex- and edge-sets are denoted by $V(M)$ and $E(M)$, respectively. If $v_1, v_2 \in V(M)$ are adjacent in $M$ we say they are $M$-\textit{adjacent}. A subset $X \subseteq V(M)$ is $M$-\textit{stable} if no two elements of $X$ are $M$-adjacent, and $X \subseteq V(M)$ is $M$-\textit{rich} if no infinite subset of $X$ is $M$-stable.

\stmt
(2.1) \textit{Let $T_1, T_2, \dots$ be a countable sequence of disjoint trees. Let $M$ be an infinite graph with $V(M) = V(T_1 \cup T_2 \cup \dots)$, such that for $i' > i \ge 1$, if $u \in V(T_i)$ is $M$-adjacent to $w \in V(T_{i'})$ and $v \in V(T_{i'}) - \{o(T_{i'})\}$ lies on the path of $T_{i'}$ from $o(T_{i'})$ to $w$ then $u$ is $M$-adjacent to $v$. Let $\{o(T_1), o(T_2), \dots\}$ be $M$-stable. Then there is an infinite $M$-stable set $X \subseteq V(M)$ such that $|X \cap V(T_i)| \le 1$ for each $i \ge 1$ and such that the set of heads of all edges of $T_1 \cup T_2 \cup \dots$ with tails in $X$ is $M$-rich.}

\textit{Proof.} We proceed by a variation on Nash--Williams' [7] ``minimal bad sequence'' argument. Let us say a sequence $z_1, z_2, \dots$ of elements of $V(M)$ is \textit{increasing} if $i' > i$ for all $j' > j \ge 1$ where $z_j \in V(T_i)$ and $z_{j'} \in V(T_{i'})$. Let $V_0 = V(M) - \{o(T_1), o(T_2), \dots\}$. A section is a countable increasing
